\section{Design}
\label{sec:design}

\subsection{Overview}

\sysname{} sits between the model runtime and the attention kernel. For each
attention layer it receives the query and key tiles, produces a block mask,
and hands the mask to a scheduler that issues only the surviving fetches.
Figure~\ref{fig:architecture} shows the arrangement.

\begin{figure}[t]
  \centering
  \begin{tikzpicture}[
      node distance=6mm and 8mm,
      box/.style={draw,rounded corners=2pt,align=center,font=\small,
                  minimum height=8mm,inner xsep=4pt},
      flow/.style={-{Stealth[length=2mm]},thick},
    ]
    \node[box] (runtime) {Model runtime};
    \node[box,below=of runtime] (probe) {Low-rank probe\\\scriptsize rank 4 sketch};
    \node[box,below=of probe] (sched) {Block scheduler\\\scriptsize mask to fetch plan};
    \node[box,below=of sched] (kernel) {Tiled attention kernel};
    \node[box,right=14mm of sched] (guard) {Residual guard};

    \draw[flow] (runtime) -- node[right,font=\scriptsize] {$Q,K$ tiles} (probe);
    \draw[flow] (probe) -- node[right,font=\scriptsize] {block mask} (sched);
    \draw[flow] (sched) -- node[right,font=\scriptsize] {fetch plan} (kernel);
    \draw[flow] (kernel.east) -- ++(6mm,0) |- (guard.east);
    \draw[flow] (guard) -- node[above,font=\scriptsize] {fallback} (sched);

    \begin{scope}[on background layer]
      \node[draw,dashed,rounded corners=3pt,fit=(probe)(sched)(guard),
            inner sep=3mm,label={[font=\scriptsize]above right:\sysname}] {};
    \end{scope}
  \end{tikzpicture}
  \caption{\sysname{} intercepts tiles between the runtime and the attention
  kernel. The guard observes realised residuals and can restore dense
  computation for a head without restarting the request.}
  \label{fig:architecture}
\end{figure}

\subsection{Block importance from a low-rank probe}

Let $Q_i \in \mathbb{R}^{B \times d}$ and $K_j \in \mathbb{R}^{B \times d}$ be
the tiles of query block $i$ and key block $j$. We maintain a shared random
projection $\Pi \in \mathbb{R}^{d \times r}$ with $r = 4$ and compute the
sketched score
\begin{equation}
  \label{eq:probe}
  \tilde{s}_{ij} = \frac{1}{\sqrt{d}} \left\lVert Q_i \Pi \right\rVert_{F}
                   \left\lVert K_j \Pi \right\rVert_{F}
                   + \frac{1}{\sqrt{d}} \mu_i^{\top} \nu_j,
\end{equation}
where $\mu_i$ and $\nu_j$ are the tile means. The first term bounds the
magnitude of the block's inner products and the second recovers the dominant
direction, which the norm term alone discards.

\begin{proposition}
\label{prop:bound}
For any temperature $\tau > 0$, the softmax mass of block $(i,j)$ satisfies
$S_{ij} \le \exp(\tilde{s}_{ij}/\tau) / Z_i$, where $Z_i$ is the row
normaliser of query block $i$.
\end{proposition}

The proposition lets us discard blocks by mass rather than by count: given a
budget $\varepsilon$, we retain the smallest set of blocks whose complementary
upper bound sums to less than $\varepsilon$.

\subsection{Scheduling}

The mask alone does not save traffic. A naive kernel still walks the full key
array and tests the mask per block, which issues the fetch before the test
resolves. \sysname{} instead compacts the surviving block indices into a dense
fetch list and issues that list to the memory pipeline, so a skipped block is
never requested. Algorithm~\ref{alg:schedule} states the procedure.

\begin{algorithm}[t]
  \caption{Block schedule for one query block}
  \label{alg:schedule}
  \SetAlgoLined
  \KwIn{query tile $Q_i$, key tiles $\{K_j\}$, budget $\varepsilon$}
  \KwOut{ordered fetch list $F$}
  compute $\tilde{s}_{ij}$ for all $j$ using Eq.~\eqref{eq:probe}\;
  sort blocks by $\tilde{s}_{ij}$ descending into order $\pi$\;
  $m \gets 0$; $F \gets \langle\,\rangle$\;
  \ForEach{$j$ in $\pi$}{
    \If{$m \ge 1 - \varepsilon$}{\textbf{break}\;}
    append $j$ to $F$\;
    $m \gets m + \exp(\tilde{s}_{ij}/\tau)/Z_i$\;
  }
  sort $F$ ascending by address to preserve coalescing\;
  \Return{$F$}
\end{algorithm}

The final sort matters more than it appears. Fetching in importance order
scatters the reads across the key array and costs 18 percent more time than
fetching the same blocks in address order.

\subsection{Runtime guard}

Proposition~\ref{prop:bound} bounds the discarded mass, but the bound is loose
for heads whose key tiles are nearly isotropic. The guard accumulates the
realised normaliser during the kernel and compares it against the predicted
one. If the relative gap exceeds $3\varepsilon$ for a head, that head is marked
dense for the remainder of the request. In our evaluation the guard fires on
1.9 percent of heads.
